% ============================================================
%  Lecture 3 — Graph coloring, timetabling, sport scheduling
%  (still Chapter 1: Introduction to Scheduling)
% ============================================================

\section{Graph Colouring}
\label{sec:graph-colouring}

Many scheduling and timetabling problems reduce, at their core,
to assigning \emph{labels} (colours) to objects so that
conflicting objects receive different labels.  This is precisely
the \textbf{graph colouring} problem.%
\aside{Applications: register allocation,
frequency assignment, map colouring.}

\begin{definition}[Graph Colouring---two formulations]
\label{def:graph-colouring}
Given an undirected graph $G = (V, E)$:
\begin{enumerate}[nosep]
  \item \textbf{Search version.}
    Given an integer $k$, find a \emph{proper $k$-colouring}:
    an assignment $c\colon V \to \{1,\dots,k\}$ such that
    $c(u) \neq c(v)$ for every edge $\{u,v\} \in E$.
  \item \textbf{Optimisation version.}
    Find the minimum $k$ for which a proper $k$-colouring exists.
    This minimum is the \emph{chromatic number} $\chi(G)$.
\end{enumerate}
\end{definition}

The search version asks: ``can I do it with $k$ colours?''
The optimisation version asks: ``what is the fewest colours I
need?''  Both are $\NP$-hard in general (for $k \ge 3$).

% ------------------------------------------------------------
\subsection{Greedy colouring}
\label{ssec:greedy-colouring}
% ------------------------------------------------------------

The simplest approach: process vertices one by one and assign each
the \emph{smallest colour not used by its already-coloured
neighbours}.  The quality depends entirely on the \textbf{vertex
ordering}.

\paragraph{Maximum-degree ordering.}
Sort vertices by decreasing degree (number of neighbours) and
colour in that order.  The intuition: high-degree vertices are
the most constrained, so handle them first while many colours
are still available.

\begin{example}[Greedy colouring by max degree]
\label{ex:greedy-colour-maxdeg}
Consider a graph on five vertices:

\begin{figure}[H]
\centering
\begin{tikzpicture}[
    vtx/.style={circle, draw=blue!60!black, fill=blue!8,
                minimum size=8mm, font=\small\bfseries},
    >=Stealth, thick
  ]
  \node[vtx] (a) at (0, 0)   {A};
  \node[vtx] (b) at (2, 0)   {B};
  \node[vtx] (c) at (3, 1.5) {C};
  \node[vtx] (d) at (1, 2.2) {D};
  \node[vtx] (e) at (-1, 1.5){E};
  \draw (a)--(b) (a)--(d) (a)--(e)
        (b)--(c) (b)--(d)
        (c)--(d)
        (d)--(e);
\end{tikzpicture}
\caption{A small graph.  Vertex~D has the highest degree (4),
         followed by A and B (degree~3 each).}
\label{fig:greedy-colour-graph}
\end{figure}

\noindent
Degrees: $D = 4$, $A = 3$, $B = 3$, $C = 2$, $E = 2$.
Processing in order $D, A, B, C, E$:
\begin{enumerate}[nosep]
  \item $D$: no coloured neighbours $\Rightarrow$ colour~1.
  \item $A$: neighbour $D$ has colour~1 $\Rightarrow$ colour~2.
  \item $B$: neighbours $D$(1), $A$(2) $\Rightarrow$ colour~3.
  \item $C$: neighbours $B$(3), $D$(1) $\Rightarrow$ colour~2.
  \item $E$: neighbours $A$(2), $D$(1) $\Rightarrow$ colour~3.
\end{enumerate}
Result: 3 colours.  In fact $\chi(G) = 3$ here, so the greedy
found the optimum---but this is not guaranteed in general.
\end{example}

% ------------------------------------------------------------
\subsection{DSatur}
\label{ssec:dsatur}
% ------------------------------------------------------------

A smarter ordering is \textbf{DSatur}
(Degree of Saturation),%
\aside{Br\'{e}laz, 1979.  Still competitive
with far more complex heuristics.}
which chooses the next vertex \emph{dynamically} rather than
fixing the order upfront.

\begin{definition}[Saturation degree]
\label{def:saturation-degree}
The \emph{saturation degree} of a vertex~$v$ is the number of
\emph{distinct colours} already assigned to its neighbours.
\end{definition}

\begin{definition}[DSatur algorithm]
\label{def:dsatur}
Repeat until all vertices are coloured:
\begin{enumerate}[nosep]
  \item Among uncoloured vertices, pick the one with the
        \textbf{highest saturation degree}.
        Break ties by highest (uncoloured) degree; break further
        ties arbitrarily.
  \item Assign it the smallest colour not used by its neighbours.
\end{enumerate}
\end{definition}

The intuition: a vertex with high saturation is surrounded by
many different colours, so it is the most ``urgent'' to colour
next---it has the fewest options left.  By always attending to the
most constrained vertex, DSatur avoids wasteful colour
introductions more often than a static ordering.

\begin{example}[DSatur on the same graph]
\label{ex:dsatur}
Using the graph from \cref{fig:greedy-colour-graph}:
\begin{enumerate}[nosep]
  \item All saturations are 0; pick $D$ (highest degree~4)
        $\Rightarrow$ colour~1.
  \item Saturations: $A = 1$, $B = 1$, $C = 1$, $E = 1$.
        Tie; pick $A$ (degree~3) $\Rightarrow$ colour~2.
  \item Saturations: $B = 2$, $C = 1$, $E = 2$.
        Tie between $B$ and $E$; pick $B$ (degree~3)
        $\Rightarrow$ colour~3.
  \item Saturations: $C = 2$, $E = 2$.
        Tie; pick $C$ (degree~2) $\Rightarrow$ colour~2.
  \item $E$: saturation~2 $\Rightarrow$ colour~3.
\end{enumerate}
Same result (3 colours), but on larger graphs DSatur typically
outperforms static max-degree ordering.
\end{example}

\begin{remark}
DSatur is exact for bipartite graphs (always finds $\chi = 2$)
and for cycle graphs.  For general graphs it is a heuristic:
it may use more than $\chi(G)$ colours.
\end{remark}


% ============================================================
\section{Timetabling Problems}
\label{sec:timetabling}
% ============================================================

Timetabling is the task of assigning events (exams, lectures,
shifts, matches) to time slots and resources (rooms, people,
fields) subject to a rich set of constraints.%
\aside{At their core, timetabling =
graph colouring with richer
constraints.}

We survey three major families: educational, workforce, and
sport timetabling.

% ------------------------------------------------------------
\subsection{Educational timetabling}
\label{ssec:educational-timetabling}
% ------------------------------------------------------------

\paragraph{Examination timetabling.}
Assign exams to time slots so that no student has two exams at
the same time.  The conflict graph has one vertex per exam and an
edge whenever two exams share at least one enrolled student.
The problem then reduces to graph colouring (minimise the number
of time slots, or fit within a given number).  Additional soft
constraints typically include:
\begin{itemize}[nosep]
  \item spreading each student's exams over the available days
        (avoid two exams on the same day or on consecutive days);
  \item room capacity limits;
  \item preferred time slots for large exams.
\end{itemize}

\paragraph{Course timetabling.}
Assign lectures to weekly time slots and rooms.
The \emph{curriculum-based} variant%
\aside{Central to ITC-2007, Track~3.}
defines conflicts not by individual enrolments but by
\emph{curricula}: pre-defined groups of courses that share
students.  Two courses in the same curriculum cannot be scheduled
at the same time.

\begin{definition}[Curriculum-Based Course Timetabling
  (CB-CTT, simplified)]
\label{def:cb-ctt}
Given:
\begin{itemize}[nosep]
  \item a set of courses, each with a number of weekly lectures,
        a teacher, and membership in one or more curricula;
  \item a set of rooms, each with a capacity;
  \item a set of time slots (periods $\times$ days);
  \item teacher unavailability constraints.
\end{itemize}
Hard constraints:
\begin{enumerate}[nosep]
  \item[H1.] No two lectures of courses in the same curriculum
             at the same time.
  \item[H2.] No teacher teaches two courses at the same time.
  \item[H3.] Room capacity is respected.
  \item[H4.] Each course gets exactly its required number of
             lectures.
\end{enumerate}
Soft constraints (penalised in the objective):
\begin{enumerate}[nosep]
  \item[S1.] Room stability: a course should use the same room
             across all its lectures.
  \item[S2.] Minimum working days: lectures of a course should
             be spread over a minimum number of days.
  \item[S3.] Curriculum compactness: lectures of the same
             curriculum on the same day should be in adjacent
             time slots (no ``holes'' in a student's schedule).
\end{enumerate}
\end{definition}

\paragraph{High-school timetabling.}
Similar to course timetabling, but typically with tighter
constraints: each class has a fixed weekly schedule, teachers
must not have gaps between lessons, and there may be requirements
on the distribution of subjects throughout the week
(e.g.\ no more than two hours of maths per day).

% ------------------------------------------------------------
\subsection{Workforce timetabling: nurse rostering}
\label{ssec:nurse-rostering}
% ------------------------------------------------------------

Workforce timetabling assigns employees to shifts over a planning
horizon, balancing coverage requirements with labour regulations
and personal preferences.  The most studied instance is the
\textbf{nurse rostering problem} (NRP).

\begin{definition}[Nurse Rostering Problem (simplified)]
\label{def:nurse-rostering}
Given:
\begin{itemize}[nosep]
  \item a set of \emph{nurses}, each with a \emph{contract}
        (min/max working days, complete-weekend flag) and a set
        of \emph{skills} (head nurse, regular nurse, trainee,
        \ldots);
  \item a set of \emph{shift types} (Early, Late, Night, \ldots)
        with forbidden successions (e.g.\ Night $\to$ Early is
        not allowed);
  \item for each (shift, skill, day) triple, a \emph{minimum}
        and \emph{optimal} staffing requirement;
  \item a set of nurse \emph{shift-off requests}: preferred
        days off or shifts to avoid.
\end{itemize}
A solution assigns each nurse, for each day of the planning
horizon (typically one week), either a (shift, skill) pair or a
day off.
\end{definition}

\paragraph{Constraints.}
Hard constraints must always be satisfied:
\begin{enumerate}[nosep]
  \item[H1.] \emph{Single assignment:} at most one shift per
             nurse per day.
  \item[H2.] \emph{Minimum staffing:} each (shift, skill, day)
             must meet the minimum requirement.
  \item[H3.] \emph{Shift successions:} no forbidden
             day-to-day shift sequences.
  \item[H4.] \emph{Skill matching:} a nurse can only cover a
             skill she possesses.
\end{enumerate}

\noindent
Soft constraints contribute to the objective function
(lower is better), each weighted:
\begin{enumerate}[nosep]
  \item[S1.] \emph{Optimal coverage} (weight~3):
             each nurse below the optimal requirement is penalised.
  \item[S2.] \emph{Preferences} (weight~1):
             each assignment to an undesired shift is penalised.
  \item[S3.] \emph{Complete weekend} (weight~3):
             if flagged, the nurse should work both Sat and Sun or
             neither.
  \item[S4.] \emph{Total assignments} (weight~2):
             working days outside the contract's $[\min, \max]$
             range are penalised.
\end{enumerate}

\begin{example}[NR-05: a five-nurse instance]
\label{ex:nr-05}
The smallest benchmark instance has 5~nurses, 2~skills, 3~shift
types, and a one-week horizon.  The instance file follows a
structured text format:

\begin{lstlisting}[language={}, numbers=none,
  caption={Excerpt from the NR-05 instance file.}]
SKILLS = 2
HeadNurse
Nurse

SHIFT_TYPES = 3
Early 0
Late 1 Early
Night 2 Early Late

CONTRACTS = 2
FullTime (4,5) 1
PartTime (2,3) 1

NURSES = 5
Patrick FullTime 2 HeadNurse Nurse
Andrea  FullTime 2 HeadNurse Nurse
Stefaan PartTime 2 HeadNurse Nurse
Sara    PartTime 1 Nurse
Nguyen  FullTime 1 Nurse
\end{lstlisting}

\noindent
A solution is a list of (nurse, day, shift, skill) assignments.
Benchmark instances with 12, 21, 30, 40, 50, 60, and 80 nurses
are available in the course materials.%
\aside{Derived from INRC-II.}
\end{example}


% ------------------------------------------------------------
\subsection{Sport timetabling: round-robin tournaments}
\label{ssec:sport-timetabling}
% ------------------------------------------------------------

In a \emph{round-robin tournament}, every team plays every
other team exactly once (single round-robin) or twice
(double round-robin, with home and away legs).
The scheduling problem: assign matches to \emph{rounds}
(time slots) so that each team plays exactly once per round,
while respecting venue constraints (home/away balance, travel
minimisation, broadcaster requests, etc.).

\paragraph{Canonical $\mathbf{1}$-factorisation.}
For $n$ teams (with $n$ even), a single round-robin requires
$n - 1$ rounds of $n/2$ simultaneous matches.
A classical construction---the \emph{canonical pattern}---fixes
team~$n$ and rotates the remaining $n - 1$ teams around a
circle:

\begin{definition}[Canonical round-robin schedule]
\label{def:canonical-rr}
For $n$ teams ($n$ even), in round~$i$
($i = 1, \dots, n-1$):
\begin{itemize}[nosep]
  \item Team~$n$ plays team~$i$ (home/away alternates by
        parity of~$i$).
  \item For $k = 1, \dots, n/2 - 1$: team
        $((k + i - 1) \bmod (n-1)) + 1$ plays team
        $((i - k - 1) \bmod (n-1)) + 1$, with home/away
        determined by parity of~$k$.
\end{itemize}
\end{definition}

This gives a valid schedule instantly (no search needed), but
it does not optimise for travel distance or home/away patterns.
Real-world sport scheduling adds many constraints, making it
a rich combinatorial problem.

\begin{example}[Round-robin for 6 teams]
\label{ex:round-robin-6}
Applying the canonical pattern with $n = 6$:

\begin{figure}[H]
\centering
\begin{tabular}{cl}
\toprule
\textbf{Round} & \textbf{Matches} \\
\midrule
1 & 1--6, \; 3--5, \; 4--2 \\
2 & 6--2, \; 4--1, \; 5--3 \\
3 & 3--6, \; 5--2, \; 1--4 \\
4 & 6--4, \; 1--3, \; 2--5 \\
5 & 5--6, \; 2--4, \; 3--1 \\
\bottomrule
\end{tabular}
\caption{Canonical round-robin schedule for 6 teams.
         In each match ``$a$--$b$'', team~$a$ is home.
         Every team appears exactly once per round.}
\label{fig:round-robin-6}
\end{figure}
\end{example}

\begin{lstlisting}[language=C++,
  caption={Canonical round-robin pattern generator
           (compile with \texttt{g++ -Wall -o pattern pattern.cpp}).},
  label={lst:pattern-generator}]
#include <iostream>
using namespace std;

int Mod(int j, int n) {
    if (j < 0) j += n;
    j = j % n;
    return (j == 0) ? n : j;
}

int main(int argc, char* argv[]) {
    if (argc != 2) {
        cerr << "Usage: " << argv[0]
             << " <number of teams (even)>" << endl;
        return 1;
    }
    int n = atoi(argv[1]);
    if (n % 2 != 0) {
        cerr << "Number of teams must be even" << endl;
        return 1;
    }
    for (int i = 1; i < n; i++) {
        cout << "Round " << i << ": ";
        if (i % 2 == 1)
            cout << i << '-' << n << " ";
        else
            cout << n << '-' << i << " ";
        for (int k = 1; k < n / 2; k++)
            if (k % 2 == 1)
                cout << Mod(k+i, n-1) << '-'
                     << Mod(i-k, n-1) << " ";
            else
                cout << Mod(i-k, n-1) << '-'
                     << Mod(k+i, n-1) << " ";
        cout << endl;
    }
    return 0;
}
\end{lstlisting}

\begin{intermezzo}{Timetabling as extended graph colouring}
All three families of timetabling---educational, workforce, and
sport---can be viewed as generalisations of graph colouring.
In examination timetabling the connection is direct: exams are
vertices, conflicts are edges, time slots are colours.
In nurse rostering, the ``colours'' become (shift, skill) pairs
and the constraints are far richer, but the core combinatorial
structure is the same: assign labels to entities so that
incompatible assignments are avoided.  This is why graph colouring
heuristics like DSatur often serve as building blocks inside
more sophisticated timetabling solvers.
\end{intermezzo}
